\documentclass[11pt]{letter}
\usepackage[margin=1in]{geometry}
\usepackage{parskip}
\signature{Vahid Rafiei}
\address{}
\begin{document}
\begin{letter}{Editor-in-Chief\\ \emph{Water Research}}
\opening{Dear Editor,}

I am pleased to submit our manuscript, ``Watershed-scale PFAS fate and transport
across surface water and groundwater: a two-way coupled SWAT\textsuperscript{+}/MODFLOW~6
model quantifies the legacy groundwater contribution to in-stream loads,'' for
consideration as a research paper in \emph{Water Research}.

Per- and polyfluoroalkyl substances (PFAS) are now federally regulated in drinking
water, and managing them turns on \emph{where} an observed in-stream concentration
originates. Groundwater discharge is the most consequential yet least resolvable of
those pathways: a grab sample cannot separate groundwater-borne from surface-derived
PFAS, and watershed models have carried PFAS only from the surface into the subsurface,
never returning the groundwater-borne fraction to the stream --- so the groundwater
\emph{contribution} to the in-stream load, the very quantity source attribution depends
on, has remained outside the model. Our work resolves this missing direction with a
two-way coupling that carries PFAS through both compartments in one continuous mass
balance, validated against an unusually rich public record (5{,}383 water-table heads,
hundreds of groundwater PFAS observations, streambed porewater, and a multi-analyte
in-stream network) on the Rogue River, Michigan, below the Wolverine World Wide legacy
source.

The central result is that two independent lines of evidence --- a model-free
multi-analyte fingerprint decomposition and a double-counting-aware joint
surface/groundwater calibration --- converge to show that legacy groundwater discharge
is a first-order, spatially localized contributor to the lower-mainstem in-stream PFOS
load, with a systematic global sensitivity and uncertainty analysis identifying the
controls on the surface-versus-groundwater partition. Enabling this are a process-based
three-phase PFAS soil model (with air--water-interface partitioning) in
SWAT\textsuperscript{+} and PFAS groundwater transport in MODFLOW~6's native framework,
rather than a frozen RT3D coupling.

This manuscript is original, unpublished, and not under consideration elsewhere. For
transparency: three \emph{methodological} components of the broader program are reported
in separate companion manuscripts under review and cited herein --- the SWATGenX
model-generation platform, the SWAT\textsuperscript{+} engine acceleration
(\emph{Geoscientific Model Development}), and a national well-lithology inventory
(\emph{Geodata and AI}) --- but the present manuscript is the only one developing and
validating the coupled surface-water/groundwater PFAS model and its findings. There are
no competing interests, and the work received no specific external funding. All code and
data are archived openly (Zenodo and a public GitHub repository).

Thank you for considering our submission; I would be glad to suggest reviewers or
provide further information.

\closing{Sincerely,}
\end{letter}
\end{document}
